\section{Minimum Cost to Cut a Stick}
\label{sec:dp-min-cost-cut-sticks}

\problemstatement{Given a stick of length $n$ and an array
$\textit{cuts}$ of positions where cuts must be made (in any order),
where the cost of making one cut equals the \textbf{current length} of
the piece being cut, find the minimum total cost to perform all cuts.}

\begin{patternbox}
This problem is Section~\ref{sec:dp-mcm}'s interval template with the
join cost changed from a triple product to a simple length difference:
choosing which cut to make \textbf{first} among a range of pending
cuts is exactly the same kind of ``pick a split point $k$'' decision.
\end{patternbox}

\subsection*{Deriving the recurrence}

Add boundary markers $0$ and $n$ to \textit{cuts}, then sort:
$\textit{points} = [0, c_1, c_2, \ldots, c_m, n]$. Let $f(i,j)$ be the
minimum cost to make every cut strictly between $\textit{points}[i]$
and $\textit{points}[j]$ (i.e., the original cuts at sorted indices
$i{+}1, \ldots, j{-}1$), on a stick segment currently spanning
$[\textit{points}[i], \textit{points}[j]]$. Base case:
$f(i, i{+}1) = 0$ (no cuts pending between adjacent points -- nothing
to do). For $j > i+1$, try every pending cut $k$ as the \textbf{first}
one made on this segment:
\[
  f(i,j) = \min_{i < k < j} \Big( f(i,k) + f(k,j) +
  \big(\textit{points}[j] - \textit{points}[i]\big) \Big).
\]
The join cost $\textit{points}[j]-\textit{points}[i]$ is constant
across every choice of $k$ (cutting first, while the segment is still
its full original length) -- only the \emph{choice of which $k$ to cut
first} affects how the two resulting sub-problems are structured.

\subsection*{C++ implementation (memoization)}

\begin{lstlisting}
#include <bits/stdc++.h>
using namespace std;

int cutStickMemo(int i, int j, vector<int>& points, vector<vector<int>>& memo) {
    if (j - i <= 1) return 0; // no cuts pending between adjacent points
    if (memo[i][j] != -1) return memo[i][j];

    int best = INT_MAX;
    for (int k = i + 1; k < j; k++) {
        int cost = cutStickMemo(i, k, points, memo) +
                   cutStickMemo(k, j, points, memo) +
                   (points[j] - points[i]);
        best = min(best, cost);
    }
    return memo[i][j] = best;
}

int minCostToCutStick(int n, vector<int>& cuts) {
    vector<int> points = cuts;
    points.push_back(0);
    points.push_back(n);
    sort(points.begin(), points.end());

    int m = points.size();
    vector<vector<int>> memo(m, vector<int>(m, -1));
    return cutStickMemo(0, m - 1, points, memo);
}

int main() {
    vector<int> cuts = {3, 5, 1, 4};
    cout << minCostToCutStick(7, cuts) << endl; // 16
    return 0;
}
\end{lstlisting}

\subsection*{Dry run}

\dryrun{Take $n=7$, $\textit{cuts}=[3,5,1,4]$. Boundary-augmented and
sorted: $\textit{points}=[0,1,3,4,5,7]$.}

The optimal ordering cuts at position $3$ first, costing the full
stick length $7$, splitting into segment $[0,3]$ (one pending cut at
$1$) and segment $[3,7]$ (two pending cuts, at $4$ and $5$). Segment
$[0,3]$ costs $3$ for its single cut. Segment $[3,7]$ optimally cuts at
$5$ first (cost $7-3=4$), splitting into $[3,5]$ (pending cut at $4$,
cost $2$) and $[5,7]$ (no pending cuts). Total cost:
$7 \,(\text{cut at }3) + 3\,(\text{cut at }1) + 4\,(\text{cut at }5) +
2\,(\text{cut at }4) = 16$, matching $f(0,5)=16$.

\begin{complexitybox}
\textbf{Time Complexity.} $O(m^3)$ where $m$ is the number of cuts
(plus boundaries) -- $O(m^2)$ ranges, each scanning $O(m)$ split
points.
\textbf{Space Complexity.} $O(m^2)$ for the memo table.
\end{complexitybox}

\begin{takeawaybox}
The choice of \emph{which} cut to make first only matters for how the
remaining sub-problems are structured -- the cost of making any cut on
a segment of fixed current length is fixed regardless of which pending
cut is chosen, which is exactly what makes the join-cost term constant
across all choices of $k$ in this recurrence.
\end{takeawaybox}
